\section{Conclusion}
\label{sec:conclusion}

We transferred an established DoE--RSM--NBI framework from engineering and forecasting to classifier ensemble weighting,
and asked not whether it produces a Pareto front but when that front can be believed. Answering that required three arms
defined by two methodological factors --- the source of the objectives and the source of the anchors --- revalidation on
the true out-of-fold objectives, an empirical reference whose sampled core is constructed independently of the
surrogate and of every candidate method, and 4 datasets $\times$ 30 partitions.

The framework is neither uniformly safe nor uniformly unsafe: its behaviour is governed by three properties of the
problem, not by anything about the optimizer.

\emph{Model class.} Surrogate adequacy is a property of the dataset \emph{and} the metric and cannot be assumed from
in-sample fit: ROC-AUC surfaces pass an external reliability gate in 0, 5, 16 and 30 of 30 partitions, log-loss
surfaces in 30, 3, 12 and 30. Smoothness in the weights is one contributing factor and not the whole explanation. It
accounts for Santander, where the rank statistic never passes the gate and has negative external $R^{2}$ despite a
rank correlation of $0.83$ while the convex log-loss passes in every partition; it does not account for BNP Paribas,
where the smooth objective does \emph{worse} than the rank statistic (3 of 30 against 5 of 30, median external
$R^{2}$ of $-0.180$ against $+0.118$) with the worst relative external fit in the study, or for Porto Seguro, where
the two are effectively tied. These data establish that adequacy varies and must be measured on unseen compositions,
not what makes a dataset--metric pair fittable.

\emph{Geometry.} Among failures attributable to the surrogate pipeline, anchor misplacement had the largest observed
effect. Replacing surrogate anchors with real single-objective optima improved both primary endpoints in all 30
partitions of Santander and BNP Paribas, whose ROC-AUC surfaces pass the gate in 0 and 5 of 30, in 23 to 24 of 30 on
Porto Seguro, which passes in 16 of 30, and in 24 of 30 on UCI credit, which passes in all 30. That substitution
changes four things at once, and a control that injects only the three real optima into the surrogate-anchored arm
shows how the effect divides: the injected points alone close a median 80\% of the hypervolume gap on Santander and
90\% on BNP Paribas, but 49\% on Porto Seguro and 1\% on UCI credit. On the two datasets with the largest raw effect
most of what is measured is therefore the returned extreme solutions. The relocated construction --- NBI takes its
utopia point, normalization and search directions from the payoff matrix, so an anchor error moves the whole frame
while an interior error perturbs single subproblems --- is the cleanest reading on UCI credit and is positive but
small elsewhere. Removing the metamodel altogether, once the anchors are right, adds far less than the anchors did on
Santander and BNP Paribas, about the same on Porto Seguro, and three times as much on UCI credit, where a
subproblem-convergence effect of the surrogate-anchored arm rather than surrogate error dominates: solver geometry can
dominate where the surrogate is excellent. With four datasets, one tie and one reversal, that ordering is an
observation and not a tested claim. The first-order fix is in any case the cheapest available: make sure the real
single-objective optima a practitioner already computes reach the returned set.

\emph{Diagnosis.} An external gate on surface quality is worth its 100 extra evaluations and separates usable from
unusable surfaces, but it does not predict anchor failure, and on one dataset the two properties diverge: the partitions
where surrogate-anchored NBI collapses are exactly those where the parsimony rule selects a surface that fits held-out
compositions better and passes the gate more often. Surface accuracy and argmax accuracy are different properties, and
only the second governs NBI; the gate should therefore be paired with an anchor check --- simply the difference between
our two surrogate arms.

Two further results are practical. The classical synergism criterion $\beta_{ij} > |\beta_i - \beta_j|$ is satisfied by
the fitted surfaces for the leading classifier pair in every partition of all four datasets, yet the real 50/50 blend
beats its better member on only one of them, and across all pairs the surface over-predicts blending gains twenty times
and under-predicts none: on these ensembles a large interaction coefficient is a statement about the fit, not evidence
of complementarity, so a candidate pair must be confirmed on the real objective. Separately, the continuous cost a
differentiable method can optimize and the step cost deployment pays disagree about which of six candidate sets
attains the best hypervolume in 10, 23 to 24, 8 and 20 of 30 partitions across the four datasets (the single-objective
references are excluded from that pool), and several conclusions reverse between them. The size of the penalty depends
on the cost range that sets the normalization box, so what we claim is the qualitative statement: sets optimized on
the linear relaxation omit the cheapest supports the unoptimized design already contains. Any cost claim must name
which cost it means, and weight cleaning is necessary --- an inert weight just above the activity threshold produced a
spurious 250~ms difference in our own analysis before we caught it.

Tripling the replications from 10 to 30 changed no directional conclusion: it converted three undecided findings into
quantified ones, exposed two apparent phenomena as artifacts and supplied intervals. We report this because it is
evidence about how much replication such a study needs: ten partitions appear to identify the conclusions, while thirty
are needed to defend their magnitudes.

One external check bounds the framework's standing: NSGA-II, given the same three objectives and the same number of real
objective evaluations per replication, produced a better approximation than the metamodel-free arm on all four datasets
and both endpoints, by small but consistent margins, and markedly more evenly spaced revalidated fronts. Where
evaluations are cheap enough to run either, that is the better choice, and the distributed representation NBI is
classically credited with is attainable here by a different construction. It does not disturb the diagnosis above ---
NSGA-II is not part of the pipeline being diagnosed --- but it settles what that pipeline is worth against the
conventional alternative.

What we carry forward is not a method but a set of controls: validate the surface on compositions it has not seen;
compare each surrogate argmax against the real optimum and prefer the real one; revalidate every candidate on the true
objective before scoring; build a reference whose sampled core is constructed independently of the surrogate and of
every candidate method; and state which cost is being paid. Each is cheap relative to the pipeline it protects, and
each changed a conclusion that would otherwise have been reported. The framework remains a good \emph{description} of the ensemble-weight simplex: its coefficients are
stable to a few percent across thirty resamplings, and they say something interpretable about which components are weak.
Whether it is also a good \emph{optimizer} there depends on conditions that can now be stated in advance and checked
cheaply.